Yes. There is a particularly clean compression which gives much more than I expected.

Let a linear extension be
[
L=(x_1,\ldots,x_n),\qquad I_k(L)={x_1,\ldots,x_k}.
]

The adjacent-swap operators (\tau_i), swapping positions (i,i+1) exactly when those elements are incomparable, are standard; the literature also explicitly groups them into odd and even operators. 

### 1. Compress by forgetting alternate prefixes

Define
[
C_o(L)=(I_2,I_4,I_6,\ldots).
]

Equivalently, remember the unordered blocks
[
B_j={x_{2j-1},x_{2j}}
]
but forget their internal order.

For a fixed compressed state (F=C_o(L)):

* each (B_j) is fixed;
* if its two elements are comparable, its order is forced;
* if they are incomparable, either order is legal;
* choices for distinct (B_j)'s are independent.

Therefore

[
\boxed{C_o^{-1}(F)\cong Q^{d(F)}}
]

**exactly**, where (d(F)) is the number of incomparable 2-blocks.

So the compression decomposes the entire LE space into actual hypercubes.

Moreover, the edges inside these cubes are exactly
[
\tau_1,\tau_3,\tau_5,\ldots.
]

Do the shifted compression
[
C_e(L)=(I_1,I_3,I_5,\ldots)
]
and its fibers are cubes whose edges are exactly
[
\tau_2,\tau_4,\tau_6,\ldots.
]

Hence

[
\boxed{
G_{BK}=\text{odd cubical foliation}+\text{even cubical foliation}.
}
]

This follows directly from the standard (\tau_i) action; non-neighbouring swaps act on disjoint position pairs. 

---

### 2. Linear statistics behave perfectly under this compression

Write a pair-orientation linear statistic as
[
f(L)=a+\sum_{{x,y}\in I(P)}
c_{xy},1_{{x<_L y}}.
]

Inside one odd fiber (F), every orientation between different blocks is fixed. The only variables are the incomparable pairs (B_j).

Thus
[
f|_F
====

a_F+\sum_{j\in D(F)}c_{B_j}Z_j,
\qquad
Z_j\sim\operatorname{Bernoulli}(1/2)
]
with the (Z_j)'s independent.

So

[
\boxed{
\operatorname{Var}(f\mid C_o)
=============================

\frac14\sum_{j\in D(C_o)}c_{B_j}^2.
}
]

There are **no covariance terms whatsoever inside a compressed fiber**.

This is exactly what we wanted compression to accomplish: the ugly correlated linear-extension measure has become a product measure on every fiber.

---

### 3. Even better: it computes BK energy exactly

Normalize the BK chain by choosing one of the (n-1) adjacent positions uniformly.

Let (\mathcal E_o) denote energy from odd-position swaps. Because (f) is degree one on every cube,

[
\boxed{
\mathcal E_o(f)
===============

\frac{2}{n-1}
\mathbb E!\left[
\operatorname{Var}(f\mid C_o)
\right].
}
]

Similarly,
[
\boxed{
\mathcal E_e(f)
===============

\frac{2}{n-1}
\mathbb E!\left[
\operatorname{Var}(f\mid C_e)
\right].
}
]

Consequently,

[
\boxed{
\mathcal E_{BK}(f)
==================

\frac{2}{n-1}
\left(
\mathbb E\operatorname{Var}(f\mid C_o)
+
\mathbb E\operatorname{Var}(f\mid C_e)
\right).
}
\tag{*}
]

That identity is exact for every pair-orientation linear statistic.

Therefore its Rayleigh quotient is

[
\boxed{
R_{BK}(f)
=========

\frac{2}{n-1}
\frac{
\mathbb E\operatorname{Var}(f\mid C_o)
+
\mathbb E\operatorname{Var}(f\mid C_e)
}{
\operatorname{Var}(f)
}.
}
\tag{**}
]

This is, I think, the real result of the compression.

---

### 4. Operator formulation

Let
[
\Pi_o f=\mathbb E[f\mid C_o],
\qquad
\Pi_e f=\mathbb E[f\mid C_e].
]

These are orthogonal projections onto functions determined by the respective compressed prefix flags.

For a linear statistic,

[
\boxed{
(I-P_{BK})f
===========

\frac{2}{n-1}
(2I-\Pi_o-\Pi_e)f.
}
\tag{***}
]

So the allegedly important standard-representation eigenvalue has been converted from a complicated BK graph problem into a problem about **two conditional-expectation projections**.

Geometrically, this is an alternating-projections/principal-angle problem:

[
\operatorname{Ran}\Pi_o
\quad\text{vs.}\quad
\operatorname{Ran}\Pi_e.
]

The small eigenvalues of
[
2I-\Pi_o-\Pi_e
]
measure how close a nonconstant function can be to being simultaneously determined by both the even-prefix and odd-prefix compressions. In the usual two-projection decomposition, these are governed by principal angles.

So, assuming your agents' claim that the relevant standard eigenfunction is linear is correct, the spectral question is essentially transformed into

[
\boxed{
\text{How correlated can the odd-prefix and even-prefix compressed information be?}
}
]

rather than a generic expansion question.

---

### 5. Why this seems highly relevant to (1/3)-(2/3)

Equation ((**)) says the desired lower bound on the linear-statistic spectral gap is equivalent to showing

[
\mathbb E\Var(f\mid C_o)
+
\mathbb E\Var(f\mid C_e)
\gtrsim
\alpha_n\Var(f)
]

for the relevant (f).

That has a compelling combinatorial meaning:

> **Some substantial fraction of the variance of (f) must be visible by locally reversing one of the disjoint adjacent incomparable pairs in one of the two checkerboard decompositions.**

A hypothetical obstruction must therefore be a statistic that is nearly determined by **both**
[
(I_2,I_4,\ldots)
\quad\text{and}\quad
(I_1,I_3,\ldots).
]

But those two prefix systems interlace. Intuitively, simultaneous predictability from both is a very rigid condition.

So I would actually **stop looking for a more elaborate compression for the moment**. This simple checkerboard compression already gives:

[
\boxed{
\text{BK graph}
\longrightarrow
\text{two exact cube foliations}
\longrightarrow
\text{conditional variances}
\longrightarrow
\text{angle between two prefix-information spaces}.
}
]

That is a much sharper bridge between **prefix statistics, BK geometry, and cubes** than the generic partial-cube observation. The next thing I would try to prove is an inequality controlling the overlap/canonical correlation of (\operatorname{Ran}\Pi_o) and (\operatorname{Ran}\Pi_e), specifically on the standard representation. If that can be expressed purely in terms of a pair bias, it may connect directly to the (1/3)-(2/3) condition.
